\documentclass[type=report,theme=light,language=de]{semio}
\title{Earth / climate / meteorological visualizations}
\author{Semio}
\date{\today}
\begin{document}
\chapter{Earth / climate / meteorological visualizations}
\section{General}
% viz-covers: 35/weather-map-35
\begin{VizFigure}[title={Weather map}, width=80, height=40]
\SemioVizDemo{weather-map-35}
\end{VizFigure}
% viz-covers: 35/synoptic-chart
\begin{VizFigure}[title={Synoptic chart}, width=80, height=40]
\SemioVizDemo{synoptic-chart}
\end{VizFigure}
% viz-covers: 35/weather-fronts
\begin{VizFigure}[title={Weather fronts}, width=80, height=40]
\SemioVizDemo{weather-fronts}
\end{VizFigure}
% viz-covers: 35/pressure-contours
\begin{VizFigure}[title={Pressure contours}, width=80, height=40]
\SemioVizDemo{pressure-contours}
\end{VizFigure}
% viz-covers: 35/isobar-map
\begin{VizFigure}[title={Isobar map}, width=80, height=40]
\SemioVizDemo{isobar-map}
\end{VizFigure}
% viz-covers: 35/isotherm-map
\begin{VizFigure}[title={Isotherm map}, width=80, height=40]
\SemioVizDemo{isotherm-map}
\end{VizFigure}
% viz-covers: 35/wind-barb-map
\begin{VizFigure}[title={Wind-barb map}, width=80, height=40]
\SemioVizDemo{wind-barb-map}
\end{VizFigure}
% viz-covers: 35/wind-rose-35
\begin{VizFigure}[title={Wind rose}, width=80, height=40]
\SemioVizDemo{wind-rose-35}
\end{VizFigure}
% viz-covers: 35/skew-t-log-p-diagram
\begin{VizFigure}[title={Skew-T log-P diagram}, width=80, height=40]
\SemioVizDemo{skew-t-log-p-diagram}
\end{VizFigure}
% viz-covers: 35/tephigram
\begin{VizFigure}[title={Tephigram}, width=80, height=40]
\SemioVizDemo{tephigram}
\end{VizFigure}
% viz-covers: 35/meteogram
\begin{VizFigure}[title={Meteogram}, width=80, height=40]
\SemioVizDemo{meteogram}
\end{VizFigure}
% viz-covers: 35/climate-stripe
\begin{VizFigure}[title={Climate stripe}, width=80, height=40]
\SemioVizDemo{climate-stripe}
\end{VizFigure}
% viz-covers: 35/climate-anomaly-chart
\begin{VizFigure}[title={Climate anomaly chart}, width=80, height=40]
\SemioVizDemo{climate-anomaly-chart}
\end{VizFigure}
% viz-covers: 35/climograph
\begin{VizFigure}[title={Climograph}, width=80, height=40]
\SemioVizDemo{climograph}
\end{VizFigure}
% viz-covers: 35/walter-lieth-climate-diagram
\begin{VizFigure}[title={Walter–Lieth climate diagram}, width=80, height=40]
\SemioVizDemo{walter-lieth-climate-diagram}
\end{VizFigure}
% viz-covers: 35/hydrograph
\begin{VizFigure}[title={Hydrograph}, width=80, height=40]
\SemioVizDemo{hydrograph}
\end{VizFigure}
% viz-covers: 35/hyetograph
\begin{VizFigure}[title={Hyetograph}, width=80, height=40]
\SemioVizDemo{hyetograph}
\end{VizFigure}
% viz-covers: 35/river-profile
\begin{VizFigure}[title={River profile}, width=80, height=40]
\SemioVizDemo{river-profile}
\end{VizFigure}
% viz-covers: 35/geological-cross-section
\begin{VizFigure}[title={Geological cross-section}, width=80, height=40]
\SemioVizDemo{geological-cross-section}
\end{VizFigure}
% viz-covers: 35/stratigraphic-column
\begin{VizFigure}[title={Stratigraphic column}, width=80, height=40]
\SemioVizDemo{stratigraphic-column}
\end{VizFigure}
% viz-covers: 35/seismic-section
\begin{VizFigure}[title={Seismic section}, width=80, height=40]
\SemioVizDemo{seismic-section}
\end{VizFigure}
% viz-covers: 35/seismogram
\begin{VizFigure}[title={Seismogram}, width=80, height=40]
\SemioVizDemo{seismogram}
\end{VizFigure}
% viz-covers: 35/tectonic-map
\begin{VizFigure}[title={Tectonic map}, width=80, height=40]
\SemioVizDemo{tectonic-map}
\end{VizFigure}
\end{document}
